\documentclass[11pt,paper=a4,answers]{exam}
\usepackage{graphicx,lastpage}
\usepackage{upgreek}
\usepackage{censor}
\usepackage{gensymb}
\usepackage{amsmath}
\usepackage{multicol}
\usepackage{enumitem}
\usepackage{tasks}
\usepackage{adjustbox}
\usepackage{xcolor}
\usepackage{tfrupee}
\setlength{\voffset}{0.25in}
\usepackage{tabularx}
\newcolumntype{Y}{>{\centering\arraybackslash}X}
%==============================================================
% Customise Non-Essential Packages Below
% =============================================================
\usepackage[most]{tcolorbox}
\usepackage{eso-pic}
\usepackage{tikz}
\AddToShipoutPictureBG{%
\begin{tikzpicture}[remember picture,overlay]
\foreach \x in {-8,-4,0,4,8,12,16,20}
\foreach \y in {-8,-4,0,4,8,12,16,20} {
\node[opacity=0.05,rotate=45,anchor=center] at ([xshift=\x cm,yshift=\y cm]current page.center) {%
\begin{minipage}{2cm}
\centering
\includegraphics[width=2cm]{images/logo_black.png}\\
\small preppeo.com
\end{minipage}%
};
}
\end{tikzpicture}%
}
\printanswers

%==============================================================
% Customise Non-Essential Packages Above
% =============================================================
\definecolor{svgray}{rgb}{0.90, 0.90, 0.90}
\graphicspath{{images/}}

\usepackage[normalem]{ulem}
\renewcommand\ULthickness{1pt}   %%---> For changing thickness of underline
\setlength\ULdepth{3.5ex}%\maxdimen ---> For changing depth of underline
\renewcommand{\baselinestretch}{1}
\chead{
\includegraphics[scale=0.1,height=2cm ]{logo.png}
}
\pagestyle{empty}

\pagestyle{headandfoot}
\headrule
\cfoot{W: https://preppeo.com; M: +91-8130711689}

\begin{document}
%==============================================================
% Customise Question paper details below this
% =============================================================
\begin{center}
    \centering
    \renewcommand{\arraystretch}{1.5}
    \begin{tabularx}{\textwidth}{YYY}
    & \normalsize\bf{{\Large{{Quant}}}} & \\
    By Preppeo & \normalsize\bf{{sat\_lid\_001}} & SAT\\
    \normalsize{{\textbf{{Algebra}} }} & \normalsize{{\textbf{{Linear equations in one variable }} }} & \normalsize{{\textbf{{Solving Linear Equations}} }} \\
    \end{tabularx}
\end{center}
\par\noindent\rule{\textwidth}{1pt}
% =============================================================
% Customise Question paper details above this
% =============================================================

\begin{questions}
\pointsinrightmargin
\pointsdroppedatright
\marksnotpoints
\pointpoints{mark}{marks}
\pointformat{\boldmath\themarginpoints}
\bracketedpoints

% =============================================================
% Question Begin
% =============================================================


%Q: 1
% Category: Practice | Difficulty: Medium | Question Type: Multiple Choice
\question
A car rental agency charges a flat daily fee of 45 plus 0.25 per mile driven. If a customer was charged a total of 95 for a one-day rental, which equation represents the number of miles, $m$, the customer drove?
\begin{tasks}(2)
    \task $45m + 0.25 = 95$
    \task $0.25m + 45 = 95$
    \task $45.25m = 95$
    \task $0.25m - 45 = 95$
\end{tasks}

\begin{solution}
To model this real-world scenario, we must identify the constant and variable components of the total cost. 
\\[10pt]
The flat daily fee of 45 is a fixed value that does not change regardless of the distance traveled. The cost per mile is 0.25, which must be multiplied by the total number of miles driven, $m$. This creates the variable term $0.25m$. 
\\[10pt]
The total charge is the sum of these two distinct parts: the variable mileage cost and the fixed daily fee. Setting this sum equal to the final bill of 95 gives us the linear equation:
\\[10pt]
$0.25m + 45 = 95$

\vspace{15pt}
Answer: \colorbox{yellow!100}{B}
\end{solution}

%Q: 2
% Category: Practice | Difficulty: Medium | Question Type: Multiple Choice
\question
If $5x + 2 = 12$, what is the value of $15x + 6$?
\begin{tasks}(2)
    \task 24
    \task 30
    \task 36
    \task 42
\end{tasks}

\begin{solution}
Instead of solving for the variable $x$ and then substituting, we can utilize algebraic manipulation to find the value of the requested expression directly. 
\\[10pt]
Observe the relationship between the given equation, $5x + 2 = 12$, and the expression we need to find, $15x + 6$. Notice that each term in the required expression is exactly three times the corresponding term in the given equation:
\\[10pt]
$3 \times (5x) = 15x$ \\
$3 \times (2) = 6$ 
\\[10pt]
According to the properties of equality, if we multiply the entire left side of the equation by 3, we must also multiply the right side by 3 to maintain balance:
\\[10pt]
$3(5x + 2) = 3(12)$ \\
$15x + 6 = 36$

\vspace{15pt}
Answer: \colorbox{yellow!100}{C}
\end{solution}

%Q: 3
% Category: Practice | Difficulty: Medium | Question Type: Multiple Choice
\question
What value of $n$ satisfies the equation $6n + 180 = 540$?
\begin{tasks}(2)
    \task 60
    \task 90
    \task 120
    \task 360
\end{tasks}

\begin{solution}
To find the value of $n$, we isolate the variable by performing inverse operations in the reverse order of operations. 
\\[10pt]
First, we address the addition of 180 by subtracting 180 from both sides of the equation:
\\[10pt]
$6n = 540 - 180$ \\
$6n = 360$ 
\\[10pt]
Next, we isolate $n$ by undoing the multiplication by 6. We divide both sides of the equation by the coefficient 6:
\\[10pt]
$n = \frac{360}{6}$ \\
$n = 60$ 
\\[10pt]
Substituting 60 back into the original equation ($6 \times 60 + 180 = 360 + 180 = 540$) confirms the solution.

\vspace{15pt}
Answer: \colorbox{yellow!100}{A}
\end{solution}

%Q: 4
% Category: Practice | Difficulty: Medium | Question Type: Multiple Choice
\question
$12x - 5x + 140 = 280$. What value of $x$ is the solution to the given equation?
\begin{tasks}(2)
    \task 20
    \task 40
    \task 70
    \task 140
\end{tasks}

\begin{solution}
The first step in solving this linear equation is to simplify the left side by combining like terms. Like terms are terms that contain the same variable raised to the same power. 
\\[10pt]
Subtracting $5x$ from $12x$ results in $7x$:
\\[10pt]
$7x + 140 = 280$ 
\\[10pt]
Now, we isolate the variable term by subtracting 140 from both sides:
\\[10pt]
$7x = 280 - 140$ \\
$7x = 140$ 
\\[10pt]
Finally, we divide both sides by 7 to solve for $x$:
\\[10pt]
$x = \frac{140}{7} = 20$

\vspace{15pt}
Answer: \colorbox{yellow!100}{A}
\end{solution}

%Q: 5
% Category: Practice | Difficulty: Easy | Question Type: Multiple Choice
\question
$(k + 9) + 21 = 50$. What value of $k$ is the solution to the given equation?
\begin{tasks}(2)
    \task 20
    \task 29
    \task 38
    \task 41
\end{tasks}

\begin{solution}
Because the operations on the left side of the equation consist only of addition, the parentheses around the expression $(k + 9)$ can be removed without affecting the value. 
\\[10pt]
The equation simplifies to:
\\[10pt]
$k + 9 + 21 = 50$ 
\\[10pt]
Next, we combine the constant terms on the left side:
\\[10pt]
$k + 30 = 50$ 
\\[10pt]
To isolate $k$, we subtract 30 from both sides of the equation:
\\[10pt]
$k = 50 - 30$ \\
$k = 20$

\vspace{15pt}
Answer: \colorbox{yellow!100}{A}
\end{solution}

%Q: 6
% Category: Practice | Difficulty: Medium | Question Type: Multiple Choice
\question
If $4x - 24 = 16$, what is the value of $x - 6$?
\begin{tasks}(2)
    \task 4
    \task 10
    \task 12
    \task 40
\end{tasks}

\begin{solution}
SAT questions often reward students for noticing patterns that bypass solving for the variable $x$. 
\\[10pt]
Observe the left side of the equation, $4x - 24$. Both terms are divisible by 4. If we factor out 4, we can rewrite the equation as:
\\[10pt]
$4(x - 6) = 16$ 
\\[10pt]
Since the question specifically asks for the value of the expression $x - 6$, we can isolate that expression by dividing both sides of the equation by 4:
\\[10pt]
$\frac{4(x - 6)}{4} = \frac{16}{4}$ \\
$x - 6 = 4$

\vspace{15pt}
Answer: \colorbox{yellow!100}{A}
\end{solution}

%Q: 7
% Category: Practice | Difficulty: Medium | Question Type: Multiple Choice
\question
$6(2x - 5) = 42$. Which equation has the same solution?
\begin{tasks}(2)
    \task $2x - 5 = 7$
    \task $2x - 5 = 36$
    \task $12x - 5 = 42$
    \task $12x - 30 = 7$
\end{tasks}

\begin{solution}
To find an equivalent equation with the same solution, we can perform operations that preserve the equality. 
\\[10pt]
The original equation states that 6 is multiplied by the quantity $(2x - 5)$. We can undo this multiplication by dividing both sides of the equation by 6:
\\[10pt]
$\frac{6(2x - 5)}{6} = \frac{42}{6}$ \\
$2x - 5 = 7$ 
\\[10pt]
This simplified version will yield the same value for $x$ as the original equation.

\vspace{15pt}
Answer: \colorbox{yellow!100}{A}
\end{solution}

%Q: 8
% Category: Practice | Difficulty: Hard | Question Type: Multiple Choice
\question
A fuel tank contained 12,000 gallons of oil. A pump removes oil at a constant rate. After 3 hours, 10,500 gallons remain. At this rate, how many total hours will the pump have run when 7,000 gallons remain?
\begin{tasks}(2)
    \task 8
    \task 10
    \task 12
    \task 15
\end{tasks}

\begin{solution}
First, we must determine the constant rate at which the oil is being removed. In the first 3 hours, the amount of oil decreased from 12,000 to 10,500 gallons:
\\[10pt]
$12,000 - 10,500 = 1,500$ gallons removed. 
\\[10pt]
Dividing this loss by the time elapsed gives the removal rate:
\\[10pt]
$1,500 \div 3 = 500$ gallons per hour. 
\\[10pt]
Next, we calculate the total amount of oil that needs to be removed from the initial quantity to reach the target of 7,000 gallons:
\\[10pt]
$12,000 - 7,000 = 5,000$ gallons. 
\\[10pt]
Finally, we divide the total required loss by the hourly rate to find the total time:
\\[10pt]
$5,000 \div 500 = 10$ hours.

\vspace{15pt}
Answer: \colorbox{yellow!100}{B}
\end{solution}

%Q: 9
% Category: Practice | Difficulty: Medium | Question Type: Multiple Choice
\question
If $3x + 4 = 13$, what is the value of $6x + 5$?
\begin{tasks}(2)
    \task 18
    \task 23
    \task 26
    \task 31
\end{tasks}

\begin{solution}
We can find the value of $x$ first. Subtracting 4 from 13 gives:
\\[10pt]
$3x = 9 \implies x = 3$ 
\\[10pt]
Now, we substitute $x = 3$ into the second expression:
\\[10pt]
$6(3) + 5 = 18 + 5 = 23$ 
\\[10pt]
Alternatively, we can notice that $6x$ is twice $3x$. If $3x = 9$, then $6x = 18$. Adding 5 to this value results in 23.

\vspace{15pt}
Answer: \colorbox{yellow!100}{B}
\end{solution}

%Q: 10
% Category: Practice | Difficulty: Medium | Question Type: Multiple Choice
\question
$8x + 16 = 176$. What is the solution to the given equation?
\begin{tasks}(2)
    \task 20
    \task 22
    \task 24
    \task 160
\end{tasks}

\begin{solution}
To solve for $x$, we first isolate the variable term by subtracting 16 from both sides of the equation:
\\[10pt]
$8x = 176 - 16$ \\
$8x = 160$ 
\\[10pt]
Next, we divide both sides by the coefficient 8:
\\[10pt]
$x = \frac{160}{8} = 20$

\vspace{15pt}
Answer: \colorbox{yellow!100}{A}
\end{solution}

%Q: 11
% Category: Practice | Difficulty: Medium | Question Type: Multiple Choice
\question
A freelancer charged 400 for a project. This included a 100 research fee and $h$ hours of editing at 60 per hour. Which equation represents this?
\begin{tasks}(2)
    \task $100h + 60 = 400$
    \task $60h + 100 = 400$
    \task $160h = 400$
    \task $60h - 100 = 400$
\end{tasks}

\begin{solution}
The total cost of the project is made up of a one-time fixed fee and a variable cost based on time. 
\\[10pt]
The fixed research fee is 100. The variable editing cost is the product of the hourly rate (60) and the number of hours ($h$), which is $60h$. 
\\[10pt]
Combining these two parts to equal the total project cost gives:
\\[10pt]
$60h + 100 = 400$

\vspace{15pt}
Answer: \colorbox{yellow!100}{B}
\end{solution}

%Q: 12
% Category: Practice | Difficulty: Medium | Question Type: Multiple Choice
\question
If $7x - 49 = 28$, what is the value of $x - 7$?
\begin{tasks}(2)
    \task 4
    \task 7
    \task 11
    \task 28
\end{tasks}

\begin{solution}
We can find the answer efficiently by factoring the left side of the equation. Since 7 is a common factor of $7x$ and 49, we can rewrite the equation as:
\\[10pt]
$7(x - 7) = 28$ 
\\[10pt]
To isolate the expression $x - 7$, divide both sides by 7:
\\[10pt]
$x - 7 = \frac{28}{7} = 4$

\vspace{15pt}
Answer: \colorbox{yellow!100}{A}
\end{solution}

%Q: 13
% Category: Practice | Difficulty: Medium | Question Type: Multiple Choice
\question
$15x - 11x + 30 = 110$. What is the value of $x$?
\begin{tasks}(2)
    \task 20
    \task 40
    \task 80
    \task 140
\end{tasks}

\begin{solution}
First, simplify the left side of the equation by combining the like $x$ terms:
\\[10pt]
$4x + 30 = 110$ 
\\[10pt]
Subtract 30 from both sides to isolate the $4x$ term:
\\[10pt]
$4x = 80$ 
\\[10pt]
Finally, divide by 4:
\\[10pt]
$x = 20$

\vspace{15pt}
Answer: \colorbox{yellow!100}{A}
\end{solution}

%Q: 14
% Category: Practice | Difficulty: Easy | Question Type: Multiple Choice
\question
$(y - 8) + 12 = 34$. What is $y$?
\begin{tasks}(2)
    \task 30
    \task 38
    \task 42
    \task 54
    \task 62
\end{tasks}

\begin{solution}
Remove the parentheses and simplify the left side by combining the numbers $-8$ and $+12$:
\\[10pt]
$y + 4 = 34$ 
\\[10pt]
Subtract 4 from both sides:
\\[10pt]
$y = 30$

\vspace{15pt}
Answer: \colorbox{yellow!100}{A}
\end{solution}

%Q: 15
% Category: Practice | Difficulty: Medium | Question Type: Multiple Choice
\question
If $2x + 1 = 6$, what is the value of $10x + 5$?
\begin{tasks}(2)
    \task 25
    \task 30
    \task 35
    \task 40
\end{tasks}

\begin{solution}
Observe the relationship between $2x + 1$ and $10x + 5$. Each term in the second expression is five times larger than the corresponding term in the first.
\\[10pt]
$5 \times (2x + 1) = 10x + 5$ 
\\[10pt]
Multiply the entire original equation by 5:
\\[10pt]
$5(2x + 1) = 5(6)$ \\
$10x + 5 = 30$

\vspace{15pt}
Answer: \colorbox{yellow!100}{B}
\end{solution}

%Q: 16
% Category: Practice | Difficulty: Medium | Question Type: Multiple Choice
\question
$3(4x + 5) = 39$. Which has the same solution?
\begin{tasks}(2)
    \task $4x + 5 = 13$
    \task $4x + 5 = 36$
    \task $12x + 5 = 39$
    \task $12x + 15 = 13$
\end{tasks}

\begin{solution}
Divide both sides of the original equation by the factor of 3 to simplify the left side:
\\[10pt]
$\frac{3(4x + 5)}{3} = \frac{39}{3}$ \\
$4x + 5 = 13$

\vspace{15pt}
Answer: \colorbox{yellow!100}{A}
\end{solution}

%Q: 17
% Category: Practice | Difficulty: Hard | Question Type: Multiple Choice
\question
A storage unit drains water at 60 liters per hour. Initially, it had 2,000 liters. After $h$ hours, 1,400 liters remain. How many total hours have passed?
\begin{tasks}(2)
    \task 8
    \task 10
    \task 12
    \task 15
\end{tasks}

\begin{solution}
Total water removed = $2,000 - 1,400 = 600$ liters. 
\\[10pt]
Given the drain rate is 60 liters per hour, we can set up the equation:
\\[10pt]
$60h = 600$ 
\\[10pt]
Divide by 60:
\\[10pt]
$h = 10$

\vspace{15pt}
Answer: \colorbox{yellow!100}{B}
\end{solution}

%Q: 18
% Category: Practice | Difficulty: Easy | Question Type: Multiple Choice
\question
What value of $p$ satisfies $9p + 270 = 900$?
\begin{tasks}(2)
    \task 70
    \task 100
    \task 130
    \task 630
\end{tasks}

\begin{solution}
Subtract 270 from both sides:
\\[10pt]
$9p = 630$ 
\\[10pt]
Divide by 9:
\\[10pt]
$p = 70$

\vspace{15pt}
Answer: \colorbox{yellow!100}{A}
\end{solution}

%Q: 19
% Category: Practice | Difficulty: Medium | Question Type: Multiple Choice
\question
If $12x - 48 = 24$, what is the value of $3x - 12$?
\begin{tasks}(2)
    \task 6
    \task 12
    \task 18
    \task 24
\end{tasks}

\begin{solution}
Each term in the expression $3x - 12$ is one-fourth of the terms in $12x - 48$. 
\\[10pt]
Divide the entire equation by 4:
\\[10pt]
$\frac{12x - 48}{4} = \frac{24}{4}$ \\
$3x - 12 = 6$

\vspace{15pt}
Answer: \colorbox{yellow!100}{A}
\end{solution}

%Q: 20
% Category: Practice | Difficulty: Easy | Question Type: Multiple Choice
\question
$18x + 10 = 190$. What is $x$?
\begin{tasks}(2)
    \task 10
    \task 12
    \task 15
    \task 180
\end{tasks}

\begin{solution}
Subtract 10:
\\[10pt]
$18x = 180$ 
\\[10pt]
Divide by 18:
\\[10pt]
$x = 10$

\vspace{15pt}
Answer: \colorbox{yellow!100}{A}
\end{solution}

%Q: 21
% Category: Practice | Difficulty: Medium | Question Type: Multiple Choice
\question
A caterer charges 150 for delivery plus 30 per guest. If the total bill was 750, how many guests, $g$, were there?
\begin{tasks}(2)
    \task $30g + 150 = 750$
    \task $150g + 30 = 750$
    \task $180g = 750$
    \task $30g = 750$
\end{tasks}

\begin{solution}
Total Cost = (Rate $\times$ guests) + Delivery Fee \\
Total = $30g + 150$. 
\\[10pt]
Given total is 750:
\\[10pt]
$30g + 150 = 750$

\vspace{15pt}
Answer: \colorbox{yellow!100}{A}
\end{solution}

%Q: 22
% Category: Practice | Difficulty: Medium | Question Type: Multiple Choice
\question
If $4x + 1 = 13$, what is $12x - 5$?
\begin{tasks}(2)
    \task 31
    \task 34
    \task 39
    \task 44
\end{tasks}

\begin{solution}
Solve for $4x$:
\\[10pt]
$4x = 12$ 
\\[10pt]
Find $12x$ by multiplying $4x$ by 3:
\\[10pt]
$12x = 3 \times 12 = 36$ 
\\[10pt]
Subtract 5:
\\[10pt]
$36 - 5 = 31$

\vspace{15pt}
Answer: \colorbox{yellow!100}{A}
\end{solution}

%Q: 23
% Category: Practice | Difficulty: Easy | Question Type: Multiple Choice
\question
$6x + 2x + 15 = 95$. Value of $x$?
\begin{tasks}(2)
    \task 10
    \task 12
    \task 15
    \task 80
\end{tasks}

\begin{solution}
Combine terms:
\\[10pt]
$8x + 15 = 95$ 
\\[10pt]
Subtract 15:
\\[10pt]
$8x = 80$ 
\\[10pt]
Divide by 8:
\\[10pt]
$x = 10$

\vspace{15pt}
Answer: \colorbox{yellow!100}{A}
\end{solution}

%Q: 24
% Category: Practice | Difficulty: Easy | Question Type: Multiple Choice
\question
$(w + 12) - 15 = 5$. What is $w$?
\begin{tasks}(2)
    \task 8
    \task 12
    \task 15
    \task 20
\end{tasks}

\begin{solution}
Simplify left side:
\\[10pt]
$w - 3 = 5$ 
\\[10pt]
Add 3:
\\[10pt]
$w = 8$

\vspace{15pt}
Answer: \colorbox{yellow!100}{A}
\end{solution}

%Q: 25
% Category: Practice | Difficulty: Easy | Question Type: Multiple Choice
\question
$3(x - 5) = 15$. Same solution as?
\begin{tasks}(2)
    \task $x - 5 = 5$
    \task $x - 5 = 12$
    \task $3x - 5 = 15$
    \task $3x - 15 = 5$
\end{tasks}

\begin{solution}
Divide by 3:
\\[10pt]
$x - 5 = 5$

\vspace{15pt}
Answer: \colorbox{yellow!100}{A}
\end{solution}

% --- PART 2: 25 NUMERIC ENTRY QUESTIONS (CORE SKILLS) ---

%Q: 26
% Category: Practice | Difficulty: Medium | Question Type: Numeric Entry
\question
Solve for $x$: $0.2x + 4 = 12$.

\begin{solution}
Isolate the decimal term by subtracting 4:
\\[10pt]
$0.2x = 8$ 
\\[10pt]
Divide by 0.2 (which is the same as multiplying by 5):
\\[10pt]
$x = 40$

\vspace{15pt}
Answer: \colorbox{yellow!100}{40}
\end{solution}

%Q: 27
% Category: Practice | Difficulty: Medium | Question Type: Numeric Entry
\question
$4(x + 2) = 2x + 14$. What is $x$?

\begin{solution}
Distribute the 4:
\\[10pt]
$4x + 8 = 2x + 14$ 
\\[10pt]
Group $x$ terms on left:
\\[10pt]
$2x + 8 = 14$ 
\\[10pt]
Subtract 8:
\\[10pt]
$2x = 6 \implies x = 3$

\vspace{15pt}
Answer: \colorbox{yellow!100}{3}
\end{solution}

%Q: 28
% Category: Practice | Difficulty: Medium | Question Type: Numeric Entry
\question
If $x/5 + 10 = 13$, what is $x$?

\begin{solution}
Subtract 10:
\\[10pt]
$x/5 = 3$ 
\\[10pt]
Multiply by 5:
\\[10pt]
$x = 15$

\vspace{15pt}
Answer: \colorbox{yellow!100}{15}
\end{solution}

%Q: 29
% Category: Practice | Difficulty: Medium | Question Type: Numeric Entry
\question
$8x - 10 = 5x + 11$. Value of $x$?

\begin{solution}
Subtract $5x$ and add 10:
\\[10pt]
$3x = 21$ 
\\[10pt]
Divide by 3:
\\[10pt]
$x = 7$

\vspace{15pt}
Answer: \colorbox{yellow!100}{7}
\end{solution}

%Q: 30
% Category: Practice | Difficulty: Medium | Question Type: Numeric Entry
\question
$2(x - 5) + 6 = 12$. What is $x$?

\begin{solution}
Distribute:
\\[10pt]
$2x - 10 + 6 = 12$ 
\\[10pt]
Simplify:
\\[10pt]
$2x - 4 = 12$ 
\\[10pt]
Add 4:
\\[10pt]
$2x = 16 \implies x = 8$

\vspace{15pt}
Answer: \colorbox{yellow!100}{8}
\end{solution}

%Q: 31
% Category: Practice | Difficulty: Easy | Question Type: Numeric Entry
\question
If $x/4 = 8$, what is $3x$?

\begin{solution}
Solve for $x$:
\\[10pt]
$x = 32$ 
\\[10pt]
Multiply by 3:
\\[10pt]
$3 \times 32 = 96$

\vspace{15pt}
Answer: \colorbox{yellow!100}{96}
\end{solution}

%Q: 32
% Category: Practice | Difficulty: Medium | Question Type: Numeric Entry
\question
$6x + 15 = 3$. What is $x$?

\begin{solution}
Subtract 15:
\\[10pt]
$6x = -12$ 
\\[10pt]
Divide by 6:
\\[10pt]
$x = -2$

\vspace{15pt}
Answer: \colorbox{yellow!100}{-2}
\end{solution}

%Q: 33
% Category: Practice | Difficulty: Medium | Question Type: Numeric Entry
\question
Solve for $m$: $5m/2 = 15$.

\begin{solution}
Multiply by 2:
\\[10pt]
$5m = 30$ 
\\[10pt]
Divide by 5:
\\[10pt]
$m = 6$

\vspace{15pt}
Answer: \colorbox{yellow!100}{6}
\end{solution}

%Q: 34
% Category: Practice | Difficulty: Medium | Question Type: Numeric Entry
\question
$25 - x = 2x + 4$. What is $x$?

\begin{solution}
Add $x$:
\\[10pt]
$25 = 3x + 4$ 
\\[10pt]
Subtract 4:
\\[10pt]
$21 = 3x \implies x = 7$

\vspace{15pt}
Answer: \colorbox{yellow!100}{7}
\end{solution}

%Q: 35
% Category: Practice | Difficulty: Easy | Question Type: Numeric Entry
\question
If $5(x + 1) = 25$, what is $x + 1$?

\begin{solution}
Divide by 5:
\\[10pt]
$x + 1 = 5$

\vspace{15pt}
Answer: \colorbox{yellow!100}{5}
\end{solution}

%Q: 36
% Category: Practice | Difficulty: Easy | Question Type: Numeric Entry
\question
$0.4x = 10$. What is $x$?

\begin{solution}
Divide by 0.4:
\\[10pt]
$x = 25$

\vspace{15pt}
Answer: \colorbox{yellow!100}{25}
\end{solution}

%Q: 37
% Category: Practice | Difficulty: Medium | Question Type: Numeric Entry
\question
$5x + 3x - 10 = 30$. What is $x$?

\begin{solution}
Combine:
\\[10pt]
$8x - 10 = 30$ 
\\[10pt]
Add 10:
\\[10pt]
$8x = 40 \implies x = 5$

\vspace{15pt}
Answer: \colorbox{yellow!100}{5}
\end{solution}

%Q: 38
% Category: Practice | Difficulty: Medium | Question Type: Numeric Entry
\question
If $12 - 4x = 28$, what is $x$?

\begin{solution}
Subtract 12:
\\[10pt]
$-4x = 16$ 
\\[10pt]
Divide by -4:
\\[10pt]
$x = -4$

\vspace{15pt}
Answer: \colorbox{yellow!100}{-4}
\end{solution}

%Q: 39
% Category: Practice | Difficulty: Medium | Question Type: Numeric Entry
\question
$4(x - 2) = x + 7$. What is $x$?

\begin{solution}
Distribute:
\\[10pt]
$4x - 8 = x + 7$ 
\\[10pt]
Subtract $x$ and add 8:
\\[10pt]
$3x = 15 \implies x = 5$

\vspace{15pt}
Answer: \colorbox{yellow!100}{5}
\end{solution}

%Q: 40
% Category: Practice | Difficulty: Medium | Question Type: Numeric Entry
\question
$x/3 - 4 = 1$. What is $x$?

\begin{solution}
Add 4:
\\[10pt]
$x/3 = 5$ 
\\[10pt]
Multiply by 3:
\\[10pt]
$x = 15$

\vspace{15pt}
Answer: \colorbox{yellow!100}{15}
\end{solution}

%Q: 41
% Category: Practice | Difficulty: Easy | Question Type: Numeric Entry
\question
If $10x = 80$, what is $x/2$?

\begin{solution}
Solve for $x$:
\\[10pt]
$x = 8$ 
\\[10pt]
Divide by 2:
\\[10pt]
$8/2 = 4$

\vspace{15pt}
Answer: \colorbox{yellow!100}{4}
\end{solution}

%Q: 42
% Category: Practice | Difficulty: Medium | Question Type: Numeric Entry
\question
$20 - (x + 5) = 10$. What is $x$?

\begin{solution}
Distribute negative:
\\[10pt]
$20 - x - 5 = 10$ 
\\[10pt]
Simplify:
\\[10pt]
$15 - x = 10 \implies x = 5$

\vspace{15pt}
Answer: \colorbox{yellow!100}{5}
\end{solution}

%Q: 43
% Category: Practice | Difficulty: Medium | Question Type: Numeric Entry
\question
$9x + 4 = 6x + 19$. Value of $x$?

\begin{solution}
Subtract $6x$:
\\[10pt]
$3x + 4 = 19$ 
\\[10pt]
Subtract 4:
\\[10pt]
$3x = 15 \implies x = 5$

\vspace{15pt}
Answer: \colorbox{yellow!100}{5}
\end{solution}

%Q: 44
% Category: Practice | Difficulty: Medium | Question Type: Numeric Entry
\question
$5(2x - 8) = 0$. What is $x$?

\begin{solution}
The quantity in parentheses must be 0:
\\[10pt]
$2x - 8 = 0$ 
\\[10pt]
$2x = 8 \implies x = 4$

\vspace{15pt}
Answer: \colorbox{yellow!100}{4}
\end{solution}

%Q: 45
% Category: Practice | Difficulty: Easy | Question Type: Numeric Entry
\question
If $x + x + x + x + x = 25$, what is $2x$?

\begin{solution}
Combine:
\\[10pt]
$5x = 25 \implies x = 5$ 
\\[10pt]
Multiply:
\\[10pt]
$2(5) = 10$

\vspace{15pt}
Answer: \colorbox{yellow!100}{10}
\end{solution}

%Q: 46
% Category: Practice | Difficulty: Medium | Question Type: Numeric Entry
\question
Solve: $1.2x + 0.3x = 3$.

\begin{solution}
Combine:
\\[10pt]
$1.5x = 3$ 
\\[10pt]
Divide by 1.5:
\\[10pt]
$x = 2$

\vspace{15pt}
Answer: \colorbox{yellow!100}{2}
\end{solution}

%Q: 47
% Category: Practice | Difficulty: Easy | Question Type: Numeric Entry
\question
If $4x = 16$, what is $4x + 9$?

\begin{solution}
Substitute:
\\[10pt]
$16 + 9 = 25$

\vspace{15pt}
Answer: \colorbox{yellow!100}{25}
\end{solution}

%Q: 48
% Category: Practice | Difficulty: Hard | Question Type: Multiple Choice
\question
$5(x + 3) = 5x + 15$. How many solutions?
\begin{tasks}(2)
    \task Zero
    \task Exactly one
    \task Infinitely many
    \task Exactly two
\end{tasks}

\begin{solution}
Distribute:
\\[10pt]
$5x + 15 = 5x + 15$ 
\\[10pt]
Identical sides mean infinite solutions.

\vspace{15pt}
Answer: \colorbox{yellow!100}{C}
\end{solution}

%Q: 49
% Category: Practice | Difficulty: Hard | Question Type: Multiple Choice
\question
$4x + 1 = 4x + 9$. How many solutions?
\begin{tasks}(2)
    \task Zero
    \task Exactly one
    \task Infinitely many
    \task Exactly two
\end{tasks}

\begin{solution}
Subtract $4x$:
\\[10pt]
$1 = 9$ 
\\[10pt]
Impossible statement means zero solutions.

\vspace{15pt}
Answer: \colorbox{yellow!100}{A}
\end{solution}

%Q: 50
% Category: Practice | Difficulty: Medium | Question Type: Numeric Entry
\question
If $1/5 x + 2/5 x = 9$, what is $x$?

\begin{solution}
Combine:
\\[10pt]
$3/5 x = 9$ 
\\[10pt]
Multiply by 5/3:
\\[10pt]
$x = 15$

\vspace{15pt}
Answer: \colorbox{yellow!100}{15}
\end{solution}


% =============================================================
% Question End
% =============================================================

\end{questions}
\begin{center}
\rule{.5\textwidth}{1pt}
\end{center}
\end{document}
